% 398_Dichte_Selbstwechselwirkung_En_ch.tex
% Johann Pascher, 18 Nov. 2026
% |psi|^2 as energy density and the self-interaction objection

\providecommand{\BM}{\textbf{[B]}}
\providecommand{\KM}{\textbf{[K]}}
\providecommand{\SM}{\textbf{[S]}}
\providecommand{\XM}{\textbf{[X]}}
\providecommand{\QM}{\textbf{[Q]}}
\providecommand{\EM}{\textbf{[E]}}

\chapter*{$|\psi|^2$ as Energy Density and the Self-Interaction Objection}
\addcontentsline{toc}{chapter}{$|\psi|^2$ as Energy Density and the Self-Interaction Objection}

\begin{abstract}
In the FFGFT, $|\psi|^2$ is not a probability but an energy density; the Born values appear as
geometric quotients (Docs.~020, 129, 230). This reading has a famous predecessor, Schrödinger's
charge density $-e|\psi|^2$ of 1926, and an equally well-known objection: if the density were real,
its parts would have to act on one another, and this self-interaction is not observed. This document
presents Schrödinger's reading and Gao's objection with verbatim quotations and recomputes it. For the
hydrogen atom the Coulomb self-energy of the 1s density would be 17.0\,eV, and the 1S--2S spacing
would shift by 12.9\,eV, more than the transition itself, which is measured to about
$4\times10^{-15}$. A self-repelling charge cloud extended in position space is therefore excluded.
The objection, however, presupposes two things: that the density is a continuum of parts in position
space, and that each part is a source for the others. The FFGFT shares neither. Its density is the
energy density of the inner configuration of one mode (Docs.~020, 230), a mode is one source with one
time scale (Doc.~334), and the field follows the realized position (Doc.~020). In field theory this
corresponds to the normal-ordered interaction term, which vanishes in the one-particle sector; this
is shown exactly. Whether the FFGFT reading delivers this property from the structure of the carrier
itself remains open. A test follows: the FFGFT expects no Schrödinger--Newton self-gravity.
\end{abstract}

% ============================================================
\section*{Status markers}
% ============================================================
% 000_Statusmarker_Legende_En_ch.tex
% Binding legend of the status markers for the entire FFGFT corpus.
% Include with % 000_Statusmarker_Legende_En_ch.tex
% Binding legend of the status markers for the entire FFGFT corpus.
% Include with % 000_Statusmarker_Legende_En_ch.tex
% Binding legend of the status markers for the entire FFGFT corpus.
% Include with \input{000_Statusmarker_Legende_En_ch} (path relative to the wrapper).
% Do not vary the wording in individual documents; check:
% python3 2/python/Korpus_Pruefung/pruef_statusmarker.py

\begin{center}
\small
\begin{tabular}{@{}lp{0.72\linewidth}@{}}
\toprule
\textbf{Marker} & \textbf{Meaning} \\
\midrule
\textbf{[K]} & Derived from $\xi$ and numerically confirmed by a check script. \\
\textbf{[K|P$n$]} & As [K], conditional on postulate P$n$. \\
\textbf{[B]} & Algebraically proved. \\
\textbf{[L]} & In addition to [B], formally proved in the proof assistant Lean~4: the Lean kernel has checked every step (sources in \texttt{2/lean/}). Applies to the mathematical statement only, not to its physical identification. \\
\textbf{[S]} & Sketched, plausible: a justified step, not yet fully derived or checked. \\
\textbf{[SETZUNG]} & Motivated postulate (axiom), not derived. \\
\textbf{[Q]} & Taken from the literature, with source. \\
\textbf{[E]} & Established physics or mathematics. \\
\textbf{[X]} & Refuted, excluded or not viable. \\
\bottomrule
\end{tabular}
\end{center}
 (path relative to the wrapper).
% Do not vary the wording in individual documents; check:
% python3 2/python/Korpus_Pruefung/pruef_statusmarker.py

\begin{center}
\small
\begin{tabular}{@{}lp{0.72\linewidth}@{}}
\toprule
\textbf{Marker} & \textbf{Meaning} \\
\midrule
\textbf{[K]} & Derived from $\xi$ and numerically confirmed by a check script. \\
\textbf{[K|P$n$]} & As [K], conditional on postulate P$n$. \\
\textbf{[B]} & Algebraically proved. \\
\textbf{[L]} & In addition to [B], formally proved in the proof assistant Lean~4: the Lean kernel has checked every step (sources in \texttt{2/lean/}). Applies to the mathematical statement only, not to its physical identification. \\
\textbf{[S]} & Sketched, plausible: a justified step, not yet fully derived or checked. \\
\textbf{[SETZUNG]} & Motivated postulate (axiom), not derived. \\
\textbf{[Q]} & Taken from the literature, with source. \\
\textbf{[E]} & Established physics or mathematics. \\
\textbf{[X]} & Refuted, excluded or not viable. \\
\bottomrule
\end{tabular}
\end{center}
 (path relative to the wrapper).
% Do not vary the wording in individual documents; check:
% python3 2/python/Korpus_Pruefung/pruef_statusmarker.py

\begin{center}
\small
\begin{tabular}{@{}lp{0.72\linewidth}@{}}
\toprule
\textbf{Marker} & \textbf{Meaning} \\
\midrule
\textbf{[K]} & Derived from $\xi$ and numerically confirmed by a check script. \\
\textbf{[K|P$n$]} & As [K], conditional on postulate P$n$. \\
\textbf{[B]} & Algebraically proved. \\
\textbf{[L]} & In addition to [B], formally proved in the proof assistant Lean~4: the Lean kernel has checked every step (sources in \texttt{2/lean/}). Applies to the mathematical statement only, not to its physical identification. \\
\textbf{[S]} & Sketched, plausible: a justified step, not yet fully derived or checked. \\
\textbf{[SETZUNG]} & Motivated postulate (axiom), not derived. \\
\textbf{[Q]} & Taken from the literature, with source. \\
\textbf{[E]} & Established physics or mathematics. \\
\textbf{[X]} & Refuted, excluded or not viable. \\
\bottomrule
\end{tabular}
\end{center}

Check script: \texttt{2/python/Dok398\_Skripte/pruef\_398\_dichte\_selbstwechselwirkung.py}.
Verbatim quotations appear in the original language in quotation marks and are marked \QM.

% ============================================================
\section{The reading of the FFGFT}
% ============================================================

The corpus fixes the meaning of $|\psi|^2$ in the same terms at several places. The amplitude follows
the local energy density of the inner configuration of the state, not a probability (Docs.~020, 384);
$|\psi|^2$ is an energy density, and the Born values appear as geometric quotients (Docs.~129, 230).
In the cylinder representation of Doc.~230, $|\alpha|^2=(1+z)/2$ is the fraction of the axis above
the point, and $r^2=1-z^2=4|\alpha|^2|\beta|^2$ is the energy density of the radial configuration
\BM. The numerical values are the same as in quantum mechanics; what differs is what they mean.
Measurement is a deterministic transition to a pole, whose outcome is fixed by $z$ and by a quantity
$\lambda$ that is set in the interaction with the measuring arrangement.

On gravitation, Doc.~020 states: the single result is fixed in the interaction with the measuring
arrangement, and the gravitational field follows the realized position. The mass is not distributed
as an average over two positions, as semiclassical gravity assumes with the expectation value of the
energy-momentum tensor as source; this fits the experiment of Page and Geilker
\QM~\cite{PageGeilker1981}. How a mass acts gravitationally before any interaction is explicitly not
derived there.

Anyone who reads $|\psi|^2$ as a real density has to face an old objection. This document presents
it, recomputes it and examines whether it applies to the reading of the FFGFT.

% ============================================================
\section{Schrödinger's charge cloud}
% ============================================================

In 1926, in the third and fourth parts of his wave mechanics, Schrödinger interpreted the quantity
$-e|\psi|^2$ as the charge density of the electron in space \QM~\cite{Schroedinger1926}: in this
reading the electron is not a point but spread over all of space, with a density that follows the
squared modulus of the wave. He showed that this density yields the correct radiation frequencies
and intensities, and he kept to this interpretation in his report to the 1927 Solvay conference
\QM~\cite{BacciagaluppiValentini2009}. Born's interpretation as a probability density prevailed.

The reading never quite disappeared. The protective measurement of Aharonov, Anandan and Vaidman
shows that $|\psi|^2$ can be measured on a single system, not only as a frequency over many runs
\QM~\cite{AharonovVaidman1993,AAV1993}: if the state is protected during a weak, long coupling, the
pointer yields the mean of the density over a small region. Gao summarizes this by saying that the
mass and charge of a single system are distributed in space, with the densities $m|\psi|^2$ and
$Q|\psi|^2$ \QM~\cite{Gao2011}; he has also worked through the history of Schrödinger's hypothesis in
its own right \QM~\cite{Gao2017}. Carver Mead explicitly treated the wave function as a matter wave
with a charge density in his \enquote{Collective Electrodynamics} \QM~\cite{Mead2000}, and Cramer and
Mead recomputed the collapse with Schrödinger's original wave mechanics as an exchange between
emitter and absorber \QM~\cite{CramerMead2020}.

For the FFGFT, protective measurement is an argument pointing in its direction: the density is a
quantity of the single system, not only of the measurement series \SM.

% ============================================================
\section{The objection}
% ============================================================

Gao states the objection as follows: if the mass and charge density is real, that is, if the
densities at different locations exist at the same time, \enquote{there will exist gravitational and
electrostatic self-interactions of the density} \QM~\cite{Gao2011}. The Schrödinger equation of a
free system with mass $m$ and charge $Q$ would then have to carry an additional term,
\begin{equation}
  i\hbar\,\partial_t\psi=-\frac{\hbar^2}{2m}\nabla^2\psi
  +\left(kQ^2-Gm^2\right)\!\int\!\frac{|\psi(x',t)|^2}{|x-x'|}\,d^3x'\;\psi ,
  \label{eq:gao}
\end{equation}
with the Coulomb constant $k$ \QM~\cite{Gao2011}. As a measure of the strength he gives, following
Salzman, $\varepsilon^2=(4Gm^2/\hbar c)^2$ for gravitation and $\varepsilon^2=(4kQ^2/\hbar c)^2$ for
electrostatics; for a free electron the electrostatic measure is about $10^{-3}$, the gravitational
one about $4\times10^{-89}$ \QM~\cite{Gao2011,Salzman2005}. The electrostatic self-interaction would
thus have to be clearly visible, in interference experiments with electrons as well as in the energy
levels of hydrogen, where the self-potential is of the same order as the potential of the nucleus.
Since neither is observed, Gao concludes that the density can \enquote{not be real but be effective}
\QM~\cite{Gao2011}: at every instant there is only one localized particle, and the density arises as
the time average of a random, discontinuous motion of this particle.

% ============================================================
\section{Recomputation}
% ============================================================

\textbf{Strength measures.} With $kQ^2=\alpha\hbar c$ the electrostatic measure is
$(4\alpha)^2=8.5\times10^{-4}$, the gravitational one $(4Gm_e^2/\hbar c)^2=4.9\times10^{-89}$ \KM.
Gao's figures agree in order of magnitude. The ratio of the two measures is about $10^{85}$ \KM.

\textbf{Hydrogen.} For the 1s density $\rho(r)=e^{-2r}/\pi$ in atomic units the self-potential of a
spherical shell is known in closed form, and the Coulomb self-energy
\begin{equation}
  J_{1s}=\int\!\!\int\frac{\rho(r)\rho(r')}{|r-r'|}\,d^3r\,d^3r'=\tfrac58\ \text{Hartree}=17.0\ \text{eV}
\end{equation}
is larger than the binding energy of $13.6$\,eV \KM. For the 2s density, numerical integration gives
$J_{2s}=77/512$\,Hartree$\,=4.09$\,eV \KM. The additional term in~\eqref{eq:gao} would shift the levels
in first order by these amounts, and hence the 1S--2S spacing by $12.9$\,eV. The transition itself
is $10.2$\,eV; the shift would be 127\,\% of it \KM. The 1S--2S frequency is measured to about 10\,Hz
at 2466\,THz, that is, to about $4\times10^{-15}$ \QM~\cite{Parthey2011}. A charge cloud in position
space whose parts repel one another according to~\eqref{eq:gao} is therefore excluded \XM.

\textbf{Gravitation.} The same density would give a gravitational self-energy of
$4\times10^{-42}$\,eV, smaller by the factor $Gm_e^2/(kQ^2)=2.4\times10^{-43}$ \KM. For an electron
the gravitational part is therefore unobservable. For heavy, extended bodies this is different; for
them there is the Schrödinger--Newton equation, which keeps exactly this term for gravitation
\QM~\cite{Diosi1984,Penrose1996}, together with proposals to measure it optomechanically
\QM~\cite{Yang2013,Grossardt2016}. Großardt et al.\ describe the equation as an
\enquote{experimentally testable alternative to quantum gravity, accessible at low energies}
\QM~\cite{Grossardt2016}.

% ============================================================
\section{What the objection presupposes}
% ============================================================

The objection is compelling if two assumptions hold.

\textbf{First:} the density is a continuum of parts in position space, and these parts are at
different places at the same time. That is Schrödinger's cloud in the literal sense.

\textbf{Second:} each part is the source of a field that acts on the other parts of the same system.
This is the assumption that equation~\eqref{eq:gao} writes down: the system's own density appears as
a source in its own potential. For gravitation this is the semiclassical coupling of Møller and
Rosenfeld, with the expectation value as source \EM.

If either assumption falls, the objection loses its force. Gao drops the first: the density becomes
effective, the product of a random motion. That costs determinism, because in his account the motion
is explicitly random and discontinuous \QM~\cite{Gao2011}. For the FFGFT this way out is not an
option.

% ============================================================
\section{The answer of the FFGFT}
% ============================================================

\subsection{Where the density lies}

The energy density of the FFGFT is not the density of a cloud in position space but that of the
inner configuration of one mode (Doc.~020). In the cylinder representation it is
$r^2=4|\alpha|^2|\beta|^2$, a quantity on the axis of the mode, not two separate charge parts at two
places (Doc.~230) \BM. A state $\alpha|0\rangle+\beta|1\rangle$ is, in this reading, one single
configuration with one point $(z,r,\theta)$, not one half here and one half there. The first
assumption of the objection does not apply to this density \SM.

\subsection{One mode, one source}

A mode has a single time scale $\tilde T=1/m$ (Doc.~334); superpositions of different modes lead to
incompatible time scales, while superposition within one mode is a configuration of this one mode
(Docs.~333, 334). In the reading of the FFGFT it follows that a mode is one source, not a continuum
of sources. Interaction couples different modes, not a mode with itself \SM.

In field theory exactly this statement has a precise form \EM. The charge density is an operator
$\hat\rho(x)=\hat\psi^\dagger(x)\hat\psi(x)$, and its expectation value in the one-particle state is
$|\psi(x)|^2$; the density is thus real in the sense of an expectation value on the single system,
as protective measurement shows it. The Coulomb term of field theory, however, is normal-ordered,
\begin{equation}
  \hat H_C=\tfrac12\!\int\!\!\int \frac{k\,{:}\hat\rho(x)\hat\rho(x'){:}}{|x-x'|}\,d^3x\,d^3x' ,
\end{equation}
and ${:}\hat\rho\hat\rho{:}$ contains two annihilators. Applied to any one-particle state it gives
zero, for every superposition \BM. The check script shows this on a lattice with an exact calculation
in Fock space: in the one-particle sector the expectation value of the two-body term is zero for all
basis states and for arbitrary superpositions, whereas a classical cloud of the same density would
have a positive self-energy; in the two-particle sector the interaction is present \KM. Real density
and absent self-interaction therefore do not exclude one another. They only do so if the density is
read as a classical continuum.

The well-known self-energy of the electron in quantum electrodynamics is something else: it comes
from the coupling to the radiation field, that is, to other modes, and enters mass renormalization
and the Lamb shift \EM. It is not a repulsion of the cloud against itself, and it fits the reading
\enquote{interaction couples different modes}.

\subsection{What the FFGFT still owes here}

That the normal-ordered form is the right one is set in field theory, not derived; it follows from
the requirement that the vacuum carries no energy from its own interaction. The FFGFT reads the same
structure as a consequence of a mode being one source with one time scale. A derivation of this
property from the structure of the carrier on $T^4$ is not available \SM.

The many-particle case is also open. For $N$ particles, $|\psi|^2$ lives in the $3N$-dimensional
configuration space. The FFGFT treats many-particle states in the product space of the modes
(Doc.~230); that the energy density there carries the same reading as for one mode is assumed, not
shown \SM.

\subsection{The gravitational test}

For gravitation the reading implies an expectation that can be measured. If the field follows the
realized position and a mode is not a source for itself, there is no Schrödinger--Newton term. The
FFGFT therefore expects no distortion of the spectrum by self-gravity in the proposed optomechanical
experiments \cite{Yang2013,Grossardt2016} \SM. A measured Schrödinger--Newton effect would contradict
this reading. The expectation is the same as that of ordinary quantum mechanics without a self-term,
but it is grounded, not merely adopted: it follows from the same statement about sources that
Doc.~020 makes for Page and Geilker. How a mass acts gravitationally before any interaction remains
open (Doc.~020); the experiments on gravitationally induced entanglement test this other part.

% ============================================================
\section{Comparison}
% ============================================================
\begin{center}
\small
\begin{tabular}{@{}p{3.0cm}p{4.6cm}p{3.0cm}p{3.6cm}@{}}
\toprule
\textbf{Reading} & \textbf{$|\psi|^2$ is} & \textbf{Single run} & \textbf{Self-interaction} \\
\midrule
Born & probability density & random & no question \\
Schrödinger (literal) & charge cloud in position space & continuous & present, \XM{} \\
Gao & effective density from random motion & random, discontinuous & none \\
Bohm & guiding field, particle separate & deterministic & none \\
FFGFT & energy density of the inner configuration of one mode & deterministic & none, one mode one source \SM \\
\bottomrule
\end{tabular}
\end{center}

Of the readings that give $|\psi|^2$ a real meaning on the single system, the FFGFT is the only one in
this list that is at the same time deterministic and reads the density itself as energy, not as a
guiding field or a time average. Its price is the open derivation of the property \enquote{one mode,
one source} from the carrier.

% ============================================================
\section{Overview}
% ============================================================
\begin{center}
\small
\begin{tabular}{@{}p{8.2cm}p{1.2cm}p{4.4cm}@{}}
\toprule
\textbf{Statement} & \textbf{Status} & \textbf{Basis} \\
\midrule
Schrödinger: $-e|\psi|^2$ as charge density & \QM & \cite{Schroedinger1926} \\
Protective measurement yields $|\psi|^2$ on the single system & \QM & \cite{AharonovVaidman1993,AAV1993} \\
Gao's objection: real density requires self-interaction & \QM & \cite{Gao2011} \\
Strength measures $8.5\times10^{-4}$ and $4.9\times10^{-89}$ & \KM & check script \\
$J_{1s}=5/8$\,Hartree, $J_{2s}=77/512$\,Hartree, 1S--2S shifted by 12.9\,eV & \KM & check script \\
Self-repelling charge cloud in position space & \XM & 1S--2S, \cite{Parthey2011} \\
$r^2=4|\alpha|^2|\beta|^2$ on the axis of one mode & \BM & Doc.~230 \\
${:}\hat\rho\hat\rho{:}$ vanishes in the one-particle sector & \BM & check script \\
FFGFT density does not meet the first assumption of the objection & \SM & Docs.~020, 230 \\
One mode, one source; derivation from the carrier & \SM & open \\
No Schrödinger--Newton self-gravity expected & \SM & test \\
\bottomrule
\end{tabular}
\end{center}

\begingroup\renewcommand{\chapter}[2]{}\small
\begin{thebibliography}{99}
\bibitem{Schroedinger1926} E.~Schrödinger, \textit{Quantisierung als Eigenwertproblem (Dritte und Vierte Mitteilung)}, Annalen der Physik \textbf{80} (1926) 437 and \textbf{81} (1926) 109.
\bibitem{BacciagaluppiValentini2009} G.~Bacciagaluppi, A.~Valentini, \textit{Quantum Theory at the Crossroads: Reconsidering the 1927 Solvay Conference}, Cambridge University Press 2009; arXiv:quant-ph/0609184.
\bibitem{AharonovVaidman1993} Y.~Aharonov, L.~Vaidman, \textit{Measurement of the Schrödinger wave of a single particle}, Physics Letters A \textbf{178} (1993) 38.
\bibitem{AAV1993} Y.~Aharonov, J.~Anandan, L.~Vaidman, \textit{Meaning of the wave function}, Physical Review A \textbf{47} (1993) 4616.
\bibitem{Gao2011} S.~Gao, \textit{The wave function and quantum reality}, arXiv:1108.1187 (2011).
\bibitem{Gao2017} S.~Gao, \textit{Is an electron a charge cloud? A reexamination of Schrödinger's charge density hypothesis}, Foundations of Science (2017), doi:10.1007/s10699-017-9521-3.
\bibitem{Salzman2005} P.~J.~Salzman, \textit{Investigation of the time dependent Schrödinger-Newton equation}, PhD thesis, University of California at Davis, 2005.
\bibitem{Mead2000} C.~A.~Mead, \textit{Collective Electrodynamics: Quantum Foundations of Electromagnetism}, MIT Press 2000.
\bibitem{CramerMead2020} J.~G.~Cramer, C.~A.~Mead, \textit{Symmetry, transactions, and the mechanism of wave function collapse}, Symmetry \textbf{12} (2020) 1373; arXiv:2006.11365.
\bibitem{Parthey2011} C.~G.~Parthey et al., \textit{Improved measurement of the hydrogen 1S--2S transition frequency}, Physical Review Letters \textbf{107} (2011) 203001.
\bibitem{Diosi1984} L.~Diósi, \textit{Gravitation and quantum-mechanical localization of macro-objects}, Physics Letters A \textbf{105} (1984) 199.
\bibitem{Penrose1996} R.~Penrose, \textit{On gravity's role in quantum state reduction}, General Relativity and Gravitation \textbf{28} (1996) 581.
\bibitem{Yang2013} H.~Yang, H.~Miao, D.-S.~Lee, B.~Helou, Y.~Chen, \textit{Macroscopic quantum mechanics in a classical spacetime}, Physical Review Letters \textbf{110} (2013) 170401.
\bibitem{Grossardt2016} A.~Großardt, J.~Bateman, H.~Ulbricht, A.~Bassi, \textit{Optomechanical test of the Schrödinger-Newton equation}, Physical Review D \textbf{93} (2016) 096003.
\bibitem{PageGeilker1981} D.~N.~Page, C.~D.~Geilker, \textit{Indirect evidence for quantum gravity}, Physical Review Letters \textbf{47} (1981) 979.
\end{thebibliography}
\endgroup
